%%%%%%%%%%%%%%%%%%%%%%%%%%%%%%%%%%%%%%%%%%%%%%%%%%%%%%%%%%%%%%%%%%%%
%% I, the copyright holder of this work, release this work into the
%% public domain. This applies worldwide. In some countries this may
%% not be legally possible; if so: I grant anyone the right to use
%% this work for any purpose, without any conditions, unless such
%% conditions are required by law.
%%%%%%%%%%%%%%%%%%%%%%%%%%%%%%%%%%%%%%%%%%%%%%%%%%%%%%%%%%%%%%%%%%%%

\documentclass{beamer}
\usetheme[faculty=ped,basePath=../fibeamer,logo=../fibeamer/logo/mu/Logo-UC-3.png]{fibeamer}
\graphicspath{{../resources/}}

\usepackage{algorithm}
\usepackage{algpseudocode}

\usetheme[faculty=ped]{fibeamer}
\usepackage[utf8]{inputenc}
\usepackage[
  main=english, %% By using `czech` or `slovak` as the main locale
                %% instead of `english`, you can typeset the
                %% presentation in either Czech or Slovak,
                %% respectively.
  czech, slovak %% The additional keys allow foreign texts to be
]{babel}        %% typeset as follows:
%%
%%   \begin{otherlanguage}{czech}   ... \end{otherlanguage}
%%   \begin{otherlanguage}{slovak}  ... \end{otherlanguage}
%%
%% These macros specify information about the presentation
\title{String} %% that will be typeset on the
\subtitle{Semestre II-2026} %% title page.
\author{Profesora: Andreina Cota Vidaurre}
%% These additional packages are used within the document:
\usepackage{ragged2e}  % `\justifying` text
\usepackage{booktabs}  % Tables
\usepackage{tabularx}
\usepackage{tikz}      % Diagrams
\usetikzlibrary{calc, shapes, backgrounds}
\usetikzlibrary{positioning, arrows.meta, shadows.blur, shapes.geometric}
\usepackage{amsmath, amssymb}
\usepackage{url}       % `\url`s
\usepackage{listings}  % Code listings
\usepackage{graphicx}
\usepackage{amsmath}
\usepackage[T1]{fontenc}
\usepackage{xcolor}
\usepackage{minted}
% \usepackage{adjustbox}


% Definicion de pseudocodigo en espaniol
\lstdefinelanguage{PseudocodigoES}{
    morekeywords={INICIO,FIN,PARA,HACER,SI,ENTONCES,FIN_SI,FIN_PARA,IMPRIMIR},
    sensitive=true,
    morecomment=[l]{//},
    morestring=[b]"
}

% Definicion de pseudocodigo en ingles
\lstdefinelanguage{PseudocodeEN}{
    morekeywords={BEGIN,END,FOR,DO,IF,THEN,END_IF,END_FOR,PRINT},
    sensitive=true,
    morecomment=[l]{//},
    morestring=[b]"
}

\lstset{
    basicstyle=\ttfamily\small,
    keywordstyle=\bfseries,
    columns=fullflexible,
    keepspaces=true,
    showstringspaces=false,
    frame=single,
    xleftmargin=0em,
    framexleftmargin=0em,
    framesep=2pt,
    gobble=4,
    resetmargins=true,
    aboveskip=0pt,
    belowskip=0pt
}

\lstdefinestyle{pythonstyle}{
    language=Python,
    basicstyle=\ttfamily\small,
    keywordstyle=\bfseries,
    commentstyle=\itshape,
    stringstyle=\ttfamily,
    showstringspaces=false,
    columns=fullflexible,
    keepspaces=true,
    frame=single,
    breaklines=true,
    xleftmargin=0em,
    framexleftmargin=0em,
    framesep=2pt,
    aboveskip=0pt,
    belowskip=0pt,
    gobble=4,
    resetmargins=true,
    emph={print,input,int,float,str,bool,len,range},
    emphstyle=\bfseries
}

\lstset{style=pythonstyle}

% Estilos del diagrama
\tikzstyle{iniciofin} = [
    ellipse,
    minimum width=2.5cm,
    minimum height=1cm,
    text centered,
    draw=black,
    fill=blue!25
]

\tikzstyle{proceso} = [
    rectangle,
    rounded corners,
    minimum width=3cm,
    minimum height=1cm,
    text centered,
    draw=black,
    fill=cyan!20
]

\tikzstyle{decision} = [
    diamond,
    aspect=2,
    text centered,
    draw=black,
    fill=yellow!35
]

\tikzstyle{flecha} = [
    thick,
    ->,
    >=Stealth
]

% \definecolor{green_python}{HTML}{52A736}

% \newcolumntype{Y}{>{\raggedright\arraybackslash}X}

\frenchspacing
\begin{document}
  \shorthandoff{-}
  \frame[c]{\maketitle}

% \begin{darkframes}

    \begin{frame}[fragile, t]{Nivel 1 -  Reconocimiento y acceso a caracteres}
        \begin{enumerate}
            \item Un soldado tiene el siguiente identificador \mintinline{python}{"BOA-1548-A"}. Mostrar:
            \begin{itemize}
                \item Cantidad de caracteres
                \item Primer caracter
                \item Último caracter
                \item Los tres primeros caracteres
                \item Los ultimos dos caracteres
            \end{itemize}
            \item Dado el código \mintinline{python}{"CAP-PEDRO-2026"}. Extraer rango militar, nombre y gestión.
            \item Dada la coordenada: \mintinline{python}{"X:153-Y:892"}. Extraer valor de X, valor de Y.
        \end{enumerate}
    \end{frame}

    \begin{frame}[fragile, t]{Nivel 2 -  Búsqueda y validación}
        \begin{enumerate}
            \item Solicitar un comando y verificar si contiene: \mintinline{python}{"AUTH"}. Si contiene la palabra, el comando es autorizado.
            \item Dado el siguiente mensaje: \mintinline{python}{"Unidad Delta reporta objetivo encontrado"}. Determinar si existe la palabra "objetivo" e indicar en qué posición comienza.
            \item Dado el mensaje: \mintinline{python}{"ALERTA ALERTA zona restringida ALERTA"}. Contar cuántas veces aparece "ALERTA".
        \end{enumerate}
    \end{frame}

    \begin{frame}[fragile, t]{Nivel 3 - Conversión y validación}
        \begin{enumerate}
            \item Solicitar el nombre de un estudiante. Luego, mostrar el nombre en mayúsculas y también en minúsculas.
            \item Solicitar un código de acceso y determinar si contiene solo números, solo letras o una mezcla.
            \item Validar la contraseña. Una contraseña es válida si tiene al menos 8 caracteres y contiene al menos una mayúscula
            y al menos un número.
        \end{enumerate}
    \end{frame}

    \begin{frame}[fragile, t]{Nivel 4 - Limpieza y transformación}
        \begin{itemize}
            \item Eliminar espacios innecesarios en el siguiente reporte: \mintinline{python}{"      Unidad lista para despliegue      "}.
            \item Dado el siguiente mensaje: \mintinline{python}{"Operacion Condor iniciada"}. Reemplazar "Condor" por "Omega".
            \item Censurar información. Dado el siguiente mensaje: \mintinline{python}{"El acceso sera por la puerta norte"}. Reemplazar \mintinline{python}{"puerta norte"} por \mintinline{python}{"*******"}.
        \end{itemize}
    \end{frame}

    \begin{frame}[fragile, t]{Nivel 5 - Reemplazar}
        \begin{itemize}
            \item En el siguiente mensaje: \mintinline{python}{"ERROR sector 1 ERROR sector 2 ERROR sector 3"}. Reemplazar solo el primer "ERROR" por "CORREGIDO". El resultado tiene que ser el siguiente: \mintinline{python}{"CORREGIDO sector 1 ERROR sector 2 ERROR sector 3"}.
            \item Sobre el mismo mensaje original, reemplazar solo el segundo "ERROR". El resultado tiene que ser el siguiente: \mintinline{python}{"ERROR sector 1 CORREGIDO sector 2 ERROR sector 3"}.
            \item Reemplazar solo la última ocurrencia de "ALERTA" por "BASE" en el siguiente mensaje: \mintinline{python}{"ALERTA base norte ALERTA base sur ALERTA central"}.
        \end{itemize}
    \end{frame}

    \begin{frame}[fragile, t]{Nivel 6 - Análisis de mensajes}
        \begin{itemize}
            \item Pedir un texto como entrada y mostrar la cantidad de vocales y la cantidad de consonantes que contiene.
            \item Verificar si un código es palíndromo. Ejemplo: \mintinline{python}{"RADAR"} o \mintinline{python}{"1221"}.
            \item Modificar el formato de los nombres de los estudiantes. Ejemplo: \mintinline{python}{"pedro perez"} por \mintinline{python}{"Pedro Perez"}.
        \end{itemize}
    \end{frame}

    \begin{frame}[fragile, t]{Nivel 7 - Problemas integradores}
        \begin{itemize}
            \item Se pide ocultar el código confidencial. Sea el texto \mintinline{python}{"4587123498765123"} modificarlo por \mintinline{python}{"************5123"}.
            \item Generar un identificador para un estudiante. Dado las siguientes variables: 
            \begin{minted}[xleftmargin=-2cm, frame=leftline, fontsize=\small]{python}
            nombre = "Carlos"
            apellido = "Mendoza"
            curso = 1001
            \end{minted}
            Generar la siguiente salida: \mintinline{python}{"cmendoza_1001"}
        \end{itemize}
    \end{frame}

    \begin{frame}[fragile,t]{Nivel 7 - Problema integradores}
        \begin{itemize}
            \item Solicitar un reporte militar y encontrar la palabra más larga.
            \item Dado el texto \mintinline{python}{"BASE-NORTE:SOLICITA-REFUERZO"}. Extraer la base y el mensaje.
            \item Dado el texto \mintinline{python}{"NORTE-145-ESTE-932"}. Extraer las coordenadas de radar.
            \item Se pide comprimir el siguiente texto \mintinline{python}{"aaabbcccc"} por \mintinline{python}{"a3b2c4"}.
        \end{itemize}
    \end{frame}    

    % \begin{frame}[fragile,t]{Nivel 8 - Concatenación y repetición}
    %     Dado un valor entero positivo N, construir una pirámide utilizando el caracter \mintinline{python}{*}. La pirámide debe tener una altura igual a N. Cada nivel de la pirámide debe estar centrado utilizando espacios en blanco.

    %     Además, cada nivel debe contener una cantidad impar de asteriscos para conservar la forma de la pirámide:
    %     \begin{itemize}
    %         \item El primer nivel contiene 1 asterisco
    %         \item El segundo nivel contiene 3 asteriscos
    %         \item El tercer nivel contiene 5 asteriscos
    %         \item Y así sucesivamente
    %     \end{itemize}
    % \end{frame}

    % \begin{frame}[fragile,t]{Nivel 8 - Concatenación y repetición}
    %     Por ejemplo:
    %     \begin{minted}[xleftmargin=-2cm, frame=leftline, fontsize=\small]{python}
    %         N = 5
    %     \end{minted}
    %     La salida esperada es: 
    %     \begin{minted}[xleftmargin=-2cm, frame=leftline, fontsize=\small]{python}
    %             *
    %            ***
    %           *****
    %          *******
    %         *********
    %     \end{minted}
    % \end{frame}

\end{document}
